% =====================================================================
%  AEGIS — MTech Thesis, IIT Delhi
%  main.tex: Master skeleton (do not write chapter content here)
% =====================================================================
\documentclass[MTech]{iitddiss}

% ------------------------------------------------------------------
% ENCODING & FONTS
% NOTE: iitddiss loads graphicx, setspace, caption automatically.
% ------------------------------------------------------------------
\usepackage[utf8]{inputenc}        % FIX: was \usepackage[utf8]{inputenc,xcolor}
\usepackage[T1]{fontenc}           % Modern form of t1enc
\usepackage{lmodern}               % Latin Modern: scalable successor to Computer Modern
                                   % eliminates bitmap .pk font warnings; same visual style
% Bold small caps: Latin Modern Roman provides no bold small-caps shape at all,
% so \textsc inside bold (e.g. \aegis in a bold heading/table cell) triggers an
% "undefined font shape" substitution warning. Map bx/sc -> bx/n (bold upright,
% which exists) so bold contexts render in bold and the warning is silenced.
% Deferred to \AtBeginDocument because T1/lmr is only registered after lmodern.
\AtBeginDocument{\DeclareFontShape{T1}{lmr}{bx}{sc}{<->ssub*lmr/bx/n}{}}
\usepackage[protrusion=true,expansion=true,final]{microtype} % full microtypography now safe
\usepackage{indentfirst}           % Indent the first line of all paragraphs (even after headings)

% ------------------------------------------------------------------
% COLOR
% ------------------------------------------------------------------
\usepackage[dvipsnames,svgnames,table]{xcolor}

% ------------------------------------------------------------------
% MATHEMATICS
% ------------------------------------------------------------------
\usepackage{amsmath}
\usepackage{amssymb}
\usepackage{amsthm}
\usepackage{mathtools}             % Extends amsmath (dcases, coloneqq, etc.)
\usepackage{bm}                    % Bold math symbols (\bm)

% ------------------------------------------------------------------
% FIGURES & GRAPHICS
% (graphicx already loaded by iitddiss.cls — no duplicate needed)
% ------------------------------------------------------------------
\usepackage{float}
\usepackage{subcaption}            % Subfigures with \subcaptionbox
\usepackage{tikz}
\usepackage{pgfplots}
\pgfplotsset{compat=1.18}
\usetikzlibrary{arrows.meta, positioning, shapes.geometric, fit, calc}
\usepackage{etoolbox}              % \patchcmd — used to inject the IITD logo into \maketitle
\graphicspath{{images/}{figures/}} % search dirs for \includegraphics

% ------------------------------------------------------------------
% TITLE-PAGE LOGO
% The title page is built by \maketitle (defined in iitddiss.cls), which
% reserves a blank 1.5in slot for the institute logo. We fill that slot
% here — from main.tex — by patching the placeholder \vspace with the
% IITD logo, keeping the title page to a single page.
%
% IMPORTANT: this patch MUST run BEFORE hyperref is loaded. hyperref
% wraps \maketitle, which moves the \vspace*{1.5in} placeholder out of
% the directly-patchable macro body and makes the patch fail.
%
% The filename contains dots (".svg.png"), so the dotted base name is
% brace-wrapped and the real extension given explicitly, per graphicx.
% ------------------------------------------------------------------
\makeatletter
\patchcmd{\maketitle}%
  {\vspace*{1.5in}}%
  {\begin{center}%
     \includegraphics[height=1.3in,keepaspectratio]%
       {{Indian_Institute_of_Technology_Delhi_Logo.svg}.png}%
   \end{center}%
   \vspace*{0.15in}}%
  {}%
  {\GenericWarning{}{main.tex: IITD logo patch into \string\maketitle\space FAILED}}%
\makeatother

% ------------------------------------------------------------------
% TABLES
% ------------------------------------------------------------------
\usepackage{booktabs}
\usepackage{multirow}
\usepackage{tabularx}
\usepackage{array}
\usepackage{siunitx}               % Aligned decimal columns in tables

% ------------------------------------------------------------------
% ALGORITHMS  (FIX: algorithmic replaced by algpseudocode)
% ------------------------------------------------------------------
\usepackage{algorithm}
\usepackage{algpseudocode}         % From algorithmicx; modern & extensible

% ------------------------------------------------------------------
% BIBLIOGRAPHY
% (iitddiss.cls calls \bibliographystyle{alpha} internally;
%  the line below overrides it to author-year, standard for EE/CS.)
% ------------------------------------------------------------------
\usepackage[numbers,sort&compress]{natbib}
\bibliographystyle{plainnat}

% ------------------------------------------------------------------
% LANGUAGE
% ------------------------------------------------------------------
\usepackage[english]{babel}

% ------------------------------------------------------------------
% SUBFILES  (lets each chapter compile independently)
% ------------------------------------------------------------------
\usepackage{subfiles}
\providecommand{\main}{.}          % Relative path to this file
\providecommand{\cwd}{.}           % Overridden by each subfile

% ------------------------------------------------------------------
% HEADER / FOOTER  (iitddiss does not set body-page headers)
% ------------------------------------------------------------------
\usepackage{fancyhdr}
\pagestyle{fancy}
\renewcommand{\sectionmark}[1]{\markright{\thesection\ #1}}
\fancyhf{}
\rhead{\fancyplain{}{\thepage}}
\lhead{\fancyplain{}{\rightmark}}
\cfoot{\textcopyright\ \the\year,\ \emph{Indian Institute of Technology Delhi}}
\renewcommand{\footrulewidth}{0.4pt}

% ------------------------------------------------------------------
% HYPERREF  — must be loaded NEAR LAST
% ------------------------------------------------------------------
% Visual options only — metadata moved to \hypersetup to avoid keyval
% comma-splitting bug (each comma in a value is misread as a new key).
\usepackage[
    colorlinks       = true,
    linkcolor        = {red!70!black},
    citecolor        = {green!50!black},
    urlcolor         = cyan,
    bookmarks        = true,
    bookmarksnumbered= true
]{hyperref}

\hypersetup{
    pdftitle    = {Aegis: A Byzantine-Resilient Federated Learning Protocol},
    pdfauthor   = {Abhishek Kumar Tripathi},
    pdfsubject  = {MTech Thesis, IIT Delhi},
    pdfkeywords = {Federated Learning, Byzantine Resilience, Aegis, Inner Product Manipulation, Sybil Attack}
}

% cleveref MUST come after hyperref
\usepackage[capitalize, noabbrev]{cleveref}

% bookmark (loaded after hyperref) rebuilds the PDF outline robustly from the
% .out file, so bookmarks appear reliably and the "main.out has changed — rerun
% to get outlines right" warning is resolved.
\usepackage{bookmark}

% ------------------------------------------------------------------
% REMOVE DOTTED LEADERS from ToC / List of Tables / List of Figures
% \@dotsep controls the spacing between the leader dots in entries
% drawn by \@dottedtocline. Setting it to a very large value pushes
% the dots infinitely apart, so none are rendered.
% ------------------------------------------------------------------
\makeatletter
\renewcommand{\@dotsep}{10000}
\makeatother

% ------------------------------------------------------------------
% THEOREM / DEFINITION ENVIRONMENTS
%
% All environments share a SINGLE running number (Theorem 1, Lemma 2,
% Assumption 3, ...). Naively sharing the [theorem] counter via
% \newtheorem{x}[theorem]{X} makes cleveref label every reference "Theorem"
% (it keys names off the counter, whose first environment is "theorem").
% thmtools' "sibling=theorem" shares the running number while "name=" registers
% the correct cleveref word per environment (Assumption, Lemma, ...). This is
% the amsthm-compatible replacement for the aliascnt trick, which clashes with
% amsthm's guarded \newcounter ("\c@lemma already defined").
% ------------------------------------------------------------------
\usepackage{thmtools}

\theoremstyle{plain}
\declaretheorem[name=Theorem, numberwithin=chapter]{theorem}
\declaretheorem[name=Lemma,       sibling=theorem]{lemma}
\declaretheorem[name=Corollary,   sibling=theorem]{corollary}
\declaretheorem[name=Proposition, sibling=theorem]{proposition}

\theoremstyle{definition}
\declaretheorem[name=Definition,  sibling=theorem]{definition}
\declaretheorem[name=Assumption,  sibling=theorem]{assumption}
\declaretheorem[name=Problem,     sibling=theorem]{problem}

\theoremstyle{remark}
\declaretheorem[name=Remark,      sibling=theorem]{remark}
\declaretheorem[name=Example,     sibling=theorem]{example}

% cleveref has no built-in names for the non-standard "assumption"/"problem"
% types (and thmtools' name= does not propagate to cleveref), so register the
% reference words explicitly. The [capitalize] option handles \cref vs \Cref.
\crefname{theorem}{Theorem}{Theorems}
\crefname{lemma}{Lemma}{Lemmas}
\crefname{corollary}{Corollary}{Corollaries}
\crefname{proposition}{Proposition}{Propositions}
\crefname{definition}{Definition}{Definitions}
\crefname{assumption}{Assumption}{Assumptions}
\crefname{problem}{Problem}{Problems}
\crefname{remark}{Remark}{Remarks}
\crefname{example}{Example}{Examples}

% ------------------------------------------------------------------
% CUSTOM MACROS — Aegis / FL domain notation
% ------------------------------------------------------------------

% System / protocol names — \textup prevents scit (small-caps italic) shape errors
% when macros appear inside \emph{} or bold headings; \texorpdfstring avoids
% PDF bookmark warnings when used in section titles.
\newcommand{\aegis}{\texorpdfstring{\textup{\textsc{Aegis}}}{Aegis}}

% Federated Learning shorthands
\newcommand{\FL}{Federated Learning}
% \textup guards against the missing small-caps-italic (scit) shape when these
% macros appear inside \emph{} or italic running headers (same fix as \aegis).
\newcommand{\ipm}{\textup{\textsc{IPM}}}               % Inner Product Manipulation
\newcommand{\fedavg}{\textup{\textsc{FedAvg}}}

% Math — sets and spaces
\newcommand{\R}{\mathbb{R}}
\newcommand{\N}{\mathbb{N}}
\newcommand{\E}{\mathbb{E}}                   % Expectation

% Math — operators
\DeclareMathOperator{\Var}{Var}
\DeclareMathOperator{\Med}{Med}               % Coordinate-wise median
\DeclareMathOperator{\clip}{clip}
\DeclareMathOperator{\sign}{sign}
\DeclareMathOperator{\diag}{diag}

% Math — norms and inner products
\newcommand{\norm}[1]{\left\lVert #1 \right\rVert}
\newcommand{\normsq}[1]{\left\lVert #1 \right\rVert^{2}}
\newcommand{\inner}[2]{\left\langle #1,\, #2 \right\rangle}
\newcommand{\abs}[1]{\left\lvert #1 \right\rvert}

% Math — complexity / O-notation
\newcommand{\Ocal}[1]{\mathcal{O}\!\left(#1\right)}

% FL system parameters  (notation matched to Appendix A convergence proof)
\newcommand{\nclients}{K}                     % Total number of clients (App. A: K)
\newcommand{\bclients}{f}                     % Number of Byzantine clients, f = |B| (App. A)
\newcommand{\globmod}{\mathbf{w}}             % Global model parameters
\newcommand{\locgrad}[1]{\mathbf{g}_{#1}}    % Client update vector (gradient delta) of client i
\newcommand{\agggrad}{\hat{\mathbf{g}}}       % Aegis aggregate (App. A: \hat{g}_t)
\newcommand{\honmean}{\bar{\mathbf{g}}_{H}}   % Honest mean update (App. A: \bar{g}_H)
\newcommand{\honvar}{V_{H}}                   % Honest gradient variance (App. A: V_H)
\newcommand{\dimmodel}{d}                     % Model parameter dimension
\newcommand{\kfilter}{k}                      % Aegis adaptive-threshold multiplier k^{(t)}

% Aegis-specific thresholds
\newcommand{\taucosine}{\tau}                 % Pass-1 cosine threshold (App. A: \tau)
\newcommand{\tauvar}{\tau_{\sigma}}           % Variance threshold

% Median references (App. A: m_1 preliminary, m_2 debiased) -- bold vectors
\newcommand{\medone}{\mathbf{m}_{1}}          % Pass-1 preliminary median
\newcommand{\medtwo}{\mathbf{m}_{2}}          % Pass-2 debiased reference median

% Convenience: Byzantine
\newcommand{\byz}{Byzantine}

% =====================================================================
%  DOCUMENT BEGIN
% =====================================================================
\begin{document}

% ------------------------------------------------------------------
% TITLE PAGE
% Fields: fill in advisor name and entry number before submission
% ------------------------------------------------------------------
\title{Aegis: A Byzantine-Resilient\\
       Federated Learning Protocol}

\author{Abhishek Kumar Tripathi}
\advisor{Prof. Harshan Jagadeesh,\\
         Indian Institute of Technology Delhi}
\entrynumber{2024EET2845}
\date{June 2026}
\department{Electrical Engineering}

\maketitle

% ------------------------------------------------------------------
% FRONT MATTER  (roman page numbers set by \maketitle above)
% ------------------------------------------------------------------
\chapter*{Certificate}
This is to certify that the thesis titled \textbf{Aegis: A Robust Byzantine-Resilient Aggregation Mechanism for Federated Learning}, submitted by \textbf{Abhishek Kumar Tripathi} to the Indian Institute of Technology Delhi, for the award of the degree of Master of Technology, is a record of bona fide work carried out by him under my supervision and guidance. The results contained in this thesis have not been submitted in part or full to any other University or Institute for the award of any degree or diploma.

\vspace{2cm}
\begin{flushright}
\textbf{Prof. Harshan Jagadeesh} \\
Department of Electrical Engineering \\
Indian Institute of Technology Delhi
\end{flushright}

\chapter*{Acknowledgements}
I would like to express my deepest gratitude to my supervisor, \textbf{Prof. Harshan Jagadeesh}, for his invaluable guidance, encouragement, and support throughout this project. His insights and feedback were crucial in shaping the research and the thesis.

I am also thankful to the Department of Electrical Engineering at IIT Delhi for providing the necessary resources and environment for this work. 

Finally, I would like to thank my family and friends for their constant support and motivation.

\subfile{\cwd/nonindex/abstract}

% ------------------------------------------------------------------
% NAVIGATIONAL LISTS  (single-spaced per IITD convention)
% ------------------------------------------------------------------
\begin{singlespace}
\tableofcontents
\thispagestyle{empty}

\listoftables
\addcontentsline{toc}{chapter}{LIST OF TABLES}

\listoffigures
\addcontentsline{toc}{chapter}{LIST OF FIGURES}
\end{singlespace}

\chapter*{List of Abbreviations}
\begin{tabular}{ll}
FL & Federated Learning \\
SGD & Stochastic Gradient Descent \\
IID & Independent and Identically Distributed \\
Non-IID & Non-Independent and Identically Distributed \\
MAD & Median Absolute Deviation \\
EMA & Exponential Moving Average \\
ALIE & A Little Is Enough \\
IPM & Inner Product Manipulation \\
CWMed & Coordinate-wise Median \\
DP & Differential Privacy \\
SMPC & Secure Multi-Party Computation \\
CNN & Convolutional Neural Network \\
Adam & Adaptive Moment Estimation \\
CIFAR & Canadian Institute for Advanced Research \\
GPU & Graphics Processing Unit \\
CPU & Central Processing Unit \\
ReLU & Rectified Linear Unit \\
GN & Group Normalization \\
MP & Max Pooling \\
DR & Detection Rate \\
Prec & Precision \\
Acc & Accuracy \\
pp & Percentage Points \\
LF & Label-Flipping \\
SF & Sign-Flipping \\
\end{tabular}

\chapter*{List of Notations}
\addcontentsline{toc}{chapter}{LIST OF NOTATIONS}

The notation below is consistent with the convergence analysis of
\cref{app:proofs}. Vectors are set in boldface (e.g.\ $\mathbf{g}_i$,
$\mathbf{w}$); scalars and scores are set in plain italic.

\vspace{1em}
\noindent\textbf{Model, Data, and Federated Objective}
\begin{tabular}{ll}
$\mathbf{w},\ \mathbf{w}^{(t)}$ & Global model parameters (at round $t$) \\
$\mathbf{w}_i^{(t)}$ & Local model weights returned by client $i$ \\
$\mathbf{x}_t$ & Model iterate in the convergence analysis (\cref{app:convergence}) \\
$\mathbf{g}_i,\ \mathbf{g}_i^{(t)}$ & Client update (gradient delta) of client $i$: $\mathbf{w}_i^{(t)}-\mathbf{w}^{(t)}$ \\
$d$ & Model parameter dimension \\
$K$ & Total number of clients \\
$\mathcal{D}_i$ & Local dataset of client $i$ \\
$n_i$ & Number of data samples at client $i$ \\
$S$ & Shards per client in the Label Sharding partition \\
$C$ & Number of classes \\
$E$ & Local training epochs per round \\
$F(\mathbf{w}),\ F_i(\mathbf{w})$ & Global objective and local loss of client $i$ \\
$F^{*}$ & Minimum value of the objective \\
$L$ & Lipschitz-smoothness constant \\
$\eta_t$ & Learning rate at round $t$ \\
\end{tabular}

\vspace{1em}
\noindent\textbf{Client Sets and Partitions}
\begin{tabular}{ll}
$\mathcal{H},\ \mathcal{B}$ & Sets of honest and Byzantine clients \\
$h=|\mathcal{H}|,\ f=|\mathcal{B}|$ & Number of honest and Byzantine clients \\
$f/K$ & Byzantine fraction \\
$\mathcal{S}^{(t)}$ & Set of clients selected in round $t$ \\
$S_1$ & Pass-1 survivor set (cosine screen) \\
$A_t$ & Approved (admitted) client set at round $t$ \\
$H_A^{(t)},\ B_A^{(t)}$ & Approved honest / Byzantine subsets: $\mathcal{H}\cap A_t$, $\mathcal{B}\cap A_t$ \\
\end{tabular}

\vspace{1em}
\noindent\textbf{Gradients and Variance}
\begin{tabular}{ll}
$\bar{\mathbf{g}}_H$ & Honest mean update: $\tfrac{1}{h}\sum_{i\in\mathcal{H}}\mathbf{g}_i$ \\
$\hat{\mathbf{g}},\ \hat{\mathbf{g}}_t$ & \aegis{} aggregate \\
$\widetilde{\mathbf{g}}_H,\ \widetilde{\mathbf{g}}_B$ & Normalised honest / Byzantine aggregates over $A_t$ \\
$\mathbf{g}_{\mathrm{byz}}$ & Byzantine (poisoned) update \\
$\mathbf{e}_t$ & \aegis{} aggregation error: $\hat{\mathbf{g}}_t-\bar{\mathbf{g}}_H$ \\
$V_H$ & Honest gradient variance: $\tfrac{1}{h}\sum_{i\in\mathcal{H}}\|\mathbf{g}_i-\bar{\mathbf{g}}_H\|^2$ \\
$\sigma_H^2$ & Uniform bound on the honest gradient variance \\
$\delta_t$ & Stochastic gradient noise: $\bar{\mathbf{g}}_H^{(t)} - \nabla f(\mathbf{x}_t)$ \\
$G$ & Bound on the honest update norm \\
$\mu,\ \sigma$ & Coordinate-wise mean / std of round updates (ALIE) \\
\end{tabular}

\vspace{1em}
\noindent\textbf{\aegis{} Scoring and Aggregation}
\begin{tabular}{ll}
$\mathbf{m}_1$ & Preliminary coordinate-wise median (Pass~1) \\
$\mathbf{m}_2$ & Debiased reference median (Pass~2) \\
$c_i$ & Cosine similarity of $\mathbf{g}_i$ with the median \\
$\tau$ & Pass-1 cosine threshold \\
$E_i$ & Euclidean distance between $\mathbf{g}_i$ and $\mathbf{m}_2$ \\
$P_i$ & Cosine penalty: $1-c_i$ \\
$R_i$ & Cross-round reputation penalty \\
$D_i$ & Credit-score denominator \\
$S_i$ & Credit score of client $i$ \\
$S_t$ & Total score of approved clients: $\sum_{i \in A_t} S_i^{(t)}$ \\
$\hat{S}_i,\ w_i^{(t)}$ & Normalised aggregation weight \\
$v_i$ & Raw client volume \\
$\widetilde{V}_i$ & Clipped client volume contribution of client $i$ \\
$M$ & Median of client volumes: $\mathrm{median}(v_j : j \in A_t)$ \\
$\alpha$ & Cosine penalty weight \\
$\lambda$ & Reputation penalty weight \\
$\gamma$ & Reputation EMA decay factor \\
$\varepsilon$ & Numerical-stability constant \\
$T^{(t)}$ & Adaptive MAD rejection threshold \\
$k^{(t)}$ & Adaptive threshold multiplier \\
$k_{\max},\ k_{\min},\ k_{\text{floor}}$ & Bounds on the adaptive multiplier \\
$\nu$ & Variance-sensitivity coefficient (Strategy~B) \\
$\rho^{(t)}$ & Credit-ramp factor \\
\end{tabular}

\vspace{1em}
\noindent\textbf{Robustness Certificate and Convergence (\cref{app:proofs})}
\begin{tabular}{ll}
$w_H,\ w_B$ & Total honest / Byzantine weight in $A_t$ ($w_H+w_B=1$) \\
$\Delta_m$ & Median contamination: $\|\mathbf{m}_2-\bar{\mathbf{g}}_H\|$ \\
$\kappa_H$ & \aegis{}-induced honest-distortion coefficient \\
$\kappa_A$ & Distortion coefficient in the robustness certificate \\
$\zeta_A$ & Residual-bias term ($\zeta_{\text{vol}}+\zeta_{\text{med}}$) \\
$r$ & Byzantine score-advantage ratio \\
$f_A$ & Approved Byzantine fraction: $|B_A|/(|H_A|+|B_A|)$ \\
$B_t$ & Upper bound on the aggregation error $\|\mathbf{e}_t\|$ \\
\end{tabular}

\vspace{1em}
\noindent\textbf{Attack Parameters}
\begin{tabular}{ll}
$s$ & Sign-flip amplification factor ($s>1$) \\
$z$ & ALIE stealth multiplier \\
$\varepsilon$ & IPM scale factor (context-distinguished from the stability constant) \\
$\pi$ & Label permutation (label-flip attack) \\
$m$ & Sybil identities per attacker \\
$f_{\text{eff}}$ & Effective Byzantine fraction under Sybil inflation \\
\end{tabular}


% ------------------------------------------------------------------
% Switch to arabic numbering for the main body
% ------------------------------------------------------------------
\cleardoublepage
\pagenumbering{arabic}

% ==================================================================
%  BODY CHAPTERS
%
%  Chapter 1 — Introduction
%  Chapter 2 — Background & Threat Model
%  Chapter 3 — The Aegis Protocol: Architecture & Design
%  Chapter 4 — Experimental Evaluation
%  Chapter 5 — Ablation Studies
%  Chapter 6 — Conclusion & Future Work
% ==================================================================

\subfile{\cwd/intro/intro.tex}
% Ch 1: Motivation, problem statement, contributions, thesis map

\subfile{\cwd/background/background.tex}
% Ch 2: FL foundations, 4-shard label-sharding Non-IID, Byzantine Tax, prior art,
%        formal threat model, 40% fraction assumption, attack formulations

\subfile{\cwd/methodology/aegis.tex}
% Ch 3: Aegis overview, Pass-1 cosine screening, Pass-2 median decontamination,
%        adaptive thresholding, O(kd) complexity, convergence analysis

\subfile{\cwd/experiments/results.tex}
% Ch 5: Testbed setup, datasets, baselines, per-attack robustness results,
%        mixed-attack scenarios

\subfile{\cwd/ablation/ablation.tex}
% Ch 6: Sensitivity to tau_cos, tau_sigma, Byzantine fraction f/n,
%        adaptive vs. fixed thresholding

\subfile{\cwd/conclusion/conclusion.tex}
% Ch 7: Summary of contributions, limitations, future directions

% ==================================================================
%  APPENDICES
% ==================================================================
\appendix

% cleveref names appendix chapters "Chapter A" by default because they still
% use the chapter counter. Retag chapter-counter labels defined from here on
% as "appendix" type so every \cref to them (including forward references from
% the body chapters) renders "Appendix A".
\crefname{appendix}{Appendix}{Appendices}
\Crefname{appendix}{Appendix}{Appendices}
\crefalias{chapter}{appendix}

\subfile{\cwd/appendices/appA_convergence.tex}
% Appendix A: Proof of Convergence Theorem and supporting lemmas

\subfile{\cwd/appendices/appB_hyperparams.tex}
% Appendix B: Full hyperparameter sweep tables and sensitivity plots

% ==================================================================
%  BIBLIOGRAPHY
% ==================================================================
\begin{singlespace}
    \bibliography{refs}
\end{singlespace}

\end{document}
